% ==========================================
% CHƯƠNG 5: KẾT LUẬN VÀ HƯỚNG PHÁT TRIỂN
% ==========================================
\chapter{KẾT LUẬN VÀ HƯỚNG PHÁT TRIỂN}
\label{Chapter5}

% ------------------------------------------------------------
\section{Kết luận và đóng góp đạt được}
\label{sec:ket_luan}
% ------------------------------------------------------------

Khóa luận này xây dựng PGA-UNet, là một mô hình phân đoạn tương tác cho ảnh X-quang xương có dùng hộp giới hạn làm câu nhắc. Mục tiêu của khóa luận là không để mô hình tự làm hết bước tìm tổn thương, mà để bác sĩ chẩn đoán ảnh khoanh sơ bộ vùng nghi ngờ trước, sau đó mô hình mới làm tiếp phần phân đoạn chi tiết hơn. Với bài toán ảnh X-quang xương, cách làm như vậy được xem là hợp lý hơn, vì tổn thương thường khó thấy rõ và cũng dễ bị nhiễu bởi các cấu trúc xung quanh.

Về mặt kỹ thuật, phần chính của khóa luận nằm ở cách đưa câu nhắc vào sâu trong mạng. PSG được dùng ở bộ mã hóa để điều biến đặc trưng theo vùng mà bác sĩ chẩn đoán ảnh đang quan tâm. CAD thì được dùng ở kết nối tắt giúp điều hướng nhánh giải mã, để tiếp tục tận dụng thông tin câu nhắc trong lúc khôi phục mặt nạ. Ngoài ra, hộp giới hạn cũng không được dùng trực tiếp dưới dạng tọa độ rời rạc, mà được đổi thành bản đồ nhiệt Gaussian dạng cao nguyên để tín hiệu không gian liên tục và mềm hơn.

Trong thực nghiệm, PGA-UNet cho kết quả cao hơn U-Net và Attention U-Net trong thiết lập so sánh giữa mô hình có câu nhắc và mô hình không có câu nhắc. Khi so sánh cùng điều kiện đầu vào với SAM-Med2D ở độ phân giải $256\times256$, PGA-UNet cũng tốt hơn trên cả hai bộ dữ liệu được khảo sát. Mô hình này cũng giữ được độ ổn định tốt hơn trong các kịch bản câu nhắc bị lệch tâm mà khóa luận đã mô phỏng. Phần nghiên cứu loại bỏ thành phần cho thấy hiệu quả của mô hình không đến từ một thành phần riêng lẻ, mà chủ yếu đến từ việc PSG, CAD và bản đồ nhiệt Gaussian đi cùng nhau.

Khi mô hình được huấn luyện lại độc lập trên FracAtlas, xu hướng kết quả chính vẫn được giữ. Nói vậy không có nghĩa là mô hình đã tổng quát hóa rộng. Ít nhất, kết quả này cho thấy thiết kế của PGA-UNet không chỉ phù hợp với một bộ dữ liệu duy nhất. Dù vậy, chỗ này vẫn chưa thể xem là bằng chứng cho khả năng tổng quát hóa liên miền theo nghĩa chặt, vì mô hình vẫn được huấn luyện riêng trên từng bộ dữ liệu trước khi đánh giá.

Khóa luận cũng làm thêm một thử nghiệm mở rộng với Gatekeeper dùng EfficientNet-B3 để xem PGA-UNet có thể đặt trong một luồng xử lý gồm cả ảnh bình thường và ảnh bệnh lý như thế nào. Chỗ này cho thấy nếu bước sàng lọc bỏ sót ảnh bệnh lý thì hiệu năng toàn luồng giảm mạnh. Vì thế, Gatekeeper chỉ nên đóng vai trò hỗ trợ sàng lọc hoặc ưu tiên ảnh nghi ngờ. Quyết định cuối cùng thì vẫn cần bác sĩ chẩn đoán ảnh xác nhận.

Về hiệu quả tính toán, PGA-UNet cũng có lợi thế rõ rệt trong môi trường thực nghiệm của khóa luận. Ở độ phân giải $256\times256$, mô hình nhỏ hơn SAM-Med2D khoảng 92 lần về số lượng tham số, 11,9 lần về FLOPs và có độ trễ suy luận thấp hơn trên cả GPU lẫn CPU. Nói ngắn gọn thì PGA-UNet hợp hơn với hướng làm gọn nhẹ và phản hồi nhanh.

Phần quan trọng nhất trong khóa luận này vẫn là mô hình PGA-UNet và cách khai thác câu nhắc trực quan thông qua PSG, CAD cùng bản đồ nhiệt Gaussian. Dù kết quả thực nghiệm nhìn chung là tốt, mô hình hiện vẫn mới dừng ở mức nghiên cứu trong điều kiện câu nhắc được kiểm soát. Vì vậy, nói đến chuyện triển khai trong môi trường lâm sàng thực tế thì còn sớm.

% ------------------------------------------------------------
\section{Hạn chế của nghiên cứu}
\label{sec:han_che}
% ------------------------------------------------------------

Bên cạnh các kết quả đạt được, nghiên cứu này vẫn còn một số hạn chế. Việc làm rõ các điểm hạn chế này là cần thiết để đánh giá đúng phạm vi và mức độ đóng góp của đề tài.

\begin{itemize}
    \item \textbf{Câu nhắc được mô phỏng từ nhãn gốc, kể cả khi ảnh có nhiều vùng tổn thương:}
    Hộp giới hạn trong thực nghiệm được tạo tự động từ mặt nạ nhãn gốc và điều chỉnh theo các kịch bản Bao trọn, Lệch tâm và Hỗn hợp; đối với ảnh có nhiều vùng tổn thương, số lượng và vị trí gần đúng của từng vùng cũng được lấy trực tiếp từ nhãn gốc (Mục~\ref{subsec:image_level_eval}, Chương~\ref{Chapter4}). Cách thiết lập này giúp kiểm soát thí nghiệm và đo đúng chất lượng phân đoạn khi đã có hộp đúng, nhưng chưa phản ánh đầy đủ sự khác biệt về vị trí, kích thước và thao tác giữa những bác sĩ chẩn đoán ảnh trong thực tế, cũng như chưa đánh giá một hệ thống phải tự phát hiện số lượng và vị trí tổn thương.

    Ngoài ra, quy tắc sinh sai lệch Lệch tâm (dịch ngẫu nhiên bằng 0,30 kích thước hộp) được dùng thống nhất trong cả huấn luyện và đánh giá. Vì vậy, kết quả về độ ổn định phản ánh khả năng thích nghi của mô hình với đúng dạng nhiễu đã học, chưa phải bằng chứng về mức độ sai lệch chưa từng thấy khi huấn luyện hoặc về hành vi khi câu nhắc không giao với tổn thương hoặc khi ảnh không có tổn thương nhưng vẫn được cung cấp hộp giới hạn, do dữ liệu huấn luyện phân đoạn hiện chỉ gồm ảnh bệnh lý với hộp luôn chứa đúng tổn thương.

    \item \textbf{Phạm vi dữ liệu còn hạn chế:}
    PGA-UNet được huấn luyện và đánh giá riêng trên BTXRD và FracAtlas. Cả hai đều là bộ dữ liệu công khai đã được chuẩn hóa; nghiên cứu chưa đánh giá trên dữ liệu X-quang thô từ cơ sở y tế, dữ liệu từ thiết bị chụp khác. Việc phân chia tập huấn luyện, xác thực và kiểm thử được thực hiện ở cấp độ ảnh (Mục~\ref{subsec:dataset}, Chương~\ref{Chapter4}); do siêu dữ liệu công khai không cung cấp mã định danh bệnh nhân đầy đủ, khóa luận chưa xác nhận được liệu có ảnh của cùng một bệnh nhân xuất hiện ở nhiều tập khác nhau hay không.

    \item \textbf{Thông tin không gian chỉ ở dạng hai chiều:}
    Ảnh X-quang là hình chiếu hai chiều của cấu trúc giải phẫu ba chiều. Hiện tượng chồng lấp giữa các cấu trúc xương có thể làm mờ ranh giới tổn thương. Bản đồ nhiệt Gaussian chỉ cung cấp vùng định hướng tương đối và không biểu diễn trực tiếp hình dạng phức tạp của tổn thương kéo dài, phân nhánh hoặc không liên tục.

    \item \textbf{Giới hạn về độ phân giải đầu vào:}
    Dù khóa luận đã dùng quy trình \textit{resize + padding} để giữ nguyên tỉ lệ ảnh và tránh méo hình học, việc chuẩn hóa ảnh về khung $512\times512$ hoặc $256\times256$ pixel vẫn có thể làm thất thoát nhiều thông tin không gian quan trọng. Do ảnh X-quang gốc thường có độ phân giải rất cao, bước thu nhỏ ảnh trước khi đệm nền vẫn có thể khiến các tổn thương kích thước nhỏ và ranh giới mờ bị mất chi tiết, từ đó làm giảm khả năng nhận diện của mô hình.

    \item \textbf{Chưa có đánh giá với bác sĩ chẩn đoán ảnh:}
    Nghiên cứu chưa đo thời gian thao tác, mức độ thuận tiện, độ biến thiên giữa bác sĩ chẩn đoán ảnh hoặc khả năng hiệu chỉnh kết quả trong một quy trình đọc ảnh thực tế. Vì vậy, hiệu quả hỗ trợ bác sĩ mới được suy luận từ kết quả kỹ thuật, chưa được xác nhận bằng nghiên cứu bác sĩ chẩn đoán ảnh.

    \item \textbf{Sai số tại mô-đun sàng lọc:}
    Gatekeeper và PGA-UNet được huấn luyện độc lập nên ảnh bệnh lý bị Gatekeeper bỏ sót sẽ không đến được bước phân đoạn trong kịch bản không có bác sĩ chẩn đoán ảnh xác nhận. Hạn chế này thể hiện rõ hơn trên FracAtlas, nơi Gatekeeper có Sensitivity thấp hơn BTXRD. Tuy nhiên, Gatekeeper chỉ là thử nghiệm mở rộng và không thuộc đóng góp kiến trúc chính của PGA-UNet.
\end{itemize}

% ------------------------------------------------------------
\section{Định hướng phát triển}
\label{sec:tuong_lai}
% ------------------------------------------------------------

Từ các hạn chế trên, khóa luận có thể đi tiếp theo một số hướng sau:

\begin{itemize}
    \item \textbf{Đánh giá với câu nhắc do bác sĩ chẩn đoán ảnh cung cấp:}
    Thực hiện thử nghiệm với bác sĩ chẩn đoán ảnh vẽ hộp giới hạn trên ảnh X-quang. Kết quả cần được đánh giá theo độ chính xác phân đoạn, thời gian thao tác, mức độ khác biệt giữa bác sĩ chẩn đoán ảnh và số lần hiệu chỉnh cần thiết. Đây là bước quan trọng để kiểm chứng độ ổn định thực tế của PGA-UNet ngoài các kịch bản câu nhắc mô phỏng.

    \item \textbf{Mở rộng đánh giá trên dữ liệu ngoài:}
    Kiểm định mô hình trên dữ liệu thu thập từ nhiều cơ sở y tế, thiết bị và điều kiện chụp khác nhau. Ngoài việc huấn luyện lại trên từng bộ dữ liệu, cần bổ sung thí nghiệm huấn luyện trên một miền và suy luận trực tiếp trên miền khác để đánh giá khả năng tổng quát hóa liên miền theo nghĩa chặt.

    \item \textbf{Xử lý ảnh có độ phân giải cao:}
    Phát triển cơ chế xử lý theo vùng quan tâm hoặc cửa sổ trượt có điều kiện theo câu nhắc, qua đó giữ lại chi tiết của tổn thương nhỏ mà không làm tăng quá mức chi phí tính toán. Các bộ mã hóa phân cấp cũng có thể được khảo sát, nhưng cần bảo đảm kiến trúc vẫn duy trì ưu điểm gọn nhẹ.

    \item \textbf{Học từ nhãn thưa:}
    Kết hợp một tập nhỏ có mặt nạ phân đoạn đầy đủ với tập lớn hơn chỉ có hộp giới hạn. PGA-UNet có thể được mở rộng theo hướng học bán giám sát hoặc sinh nhãn giả, qua đó giảm chi phí tạo chú thích đa giác ở cấp độ pixel.

    \item \textbf{Kết hợp thông tin lâm sàng:}
    Khảo sát cơ chế dung hợp đặc trưng hình ảnh với các thông tin như tuổi, giới tính, vị trí đau hoặc tiền sử bệnh. Thông tin bổ sung có thể hỗ trợ phân biệt các trường hợp khó, nhưng cần được đánh giá trên dữ liệu lâm sàng được thu thập và ẩn danh đúng quy trình.
\end{itemize}

Việc cần làm trước vẫn là kiểm tra PGA-UNet với câu nhắc do bác sĩ chẩn đoán ảnh thật cung cấp và trên dữ liệu ngoài môi trường công khai. Làm xong bước đó rồi mới nên tính tiếp đến các hướng như xử lý ảnh độ phân giải cao, học từ nhãn thưa, kết hợp thêm thông tin lâm sàng hoặc cải tiến Gatekeeper để hệ thống có khả năng dùng thực tế tốt hơn.
